% Reviewer #2, Comment 1 -- paste-ready manuscript insert.
% Drop into the evaluation section (after the existing Mode A / Mode B tables),
% renumber tables to match the surrounding manuscript, and keep the cross-refs.
% Requires \usepackage{booktabs,multirow}.

\subsection{Broader Adversarial Evaluation}
\label{sec:broader-adversarial}

The earlier evaluation exercised a fixed set of modelled attacks. To test whether
the guarantees hold beyond that set, we replaced the fixed list with a
\emph{compositional} attack space and re-ran every scenario against the same
deterministic verifier.

\paragraph{Compositional attack matrix.}
We decompose the threat model into nine independent attack primitives spanning
merchant, agent, settlement, delivery and evidence surfaces (payee substitution,
merchant revocation, amount escalation, payment-key de-authorization, asset
switch, chain transplant, task substitution, delivery forgery, and
selection injection). Each primitive is executed alone, in every ordered pair,
and in twelve curated three-step chains that model an attacker advancing through
discovery $\rightarrow$ payment $\rightarrow$ delivery or through identity
$\rightarrow$ amount $\rightarrow$ settlement. Every scenario is instantiated
under five seeds that vary the amount, asset, attacker payee, task hash and
token, so a detection is not an artefact of one fixed trace.

Across 465 executed scenarios the verifier detected every attack
(465/465, 100\%), and in every scenario the rules that failed included the rule
targeting the injected primitive (Table~\ref{tab:adversarial-matrix}). Detection
is uniform across chain depth and across attack families, which indicates the
guarantees compose rather than depending on a single fixture.

\begin{table}[t]
\centering
\small
\caption{Compositional adversarial matrix. Every scenario is executed against
the deterministic verifier; ``detected'' means at least one rule or the
selection-audit check failed.}
\label{tab:adversarial-matrix}
\begin{tabular}{@{}llrr@{}}
\toprule
Axis & Bucket & Scenarios & Detected \\
\midrule
\multirow{3}{*}{Chain depth} & single primitive & 45 & 45 (100\%) \\
                             & ordered pair     & 360 & 360 (100\%) \\
                             & three-step chain & 60 & 60 (100\%) \\
\midrule
\multirow{5}{*}{Family}      & merchant         & 210 & 210 (100\%) \\
                             & agent            & 215 & 215 (100\%) \\
                             & settlement       & 200 & 200 (100\%) \\
                             & delivery         & 215 & 215 (100\%) \\
                             & evidence         & 105 & 105 (100\%) \\
\midrule
\multicolumn{2}{@{}l}{Total} & 465 & 465 (100\%) \\
\bottomrule
\end{tabular}
\end{table}

\paragraph{Malicious merchants and compromised agents.}
The two adversarial roles the reviewer highlights are covered directly. For
malicious merchants we execute payee substitution, merchant-identity revocation,
and delivery forgery, alone and combined with a manipulated discovery step; all
are rejected. For compromised agents we assume control of the agent's
orchestration logic but \emph{not} its signing key, and execute amount
escalation, asset switching, chain switching, task substitution and payment-key
de-authorization; all are rejected, because the human-signed mandate and the
recomputed commitment bind the fields the agent would have to change.

\paragraph{Pre-settlement prevention.}
Detection after settlement is necessary but not sufficient, so we additionally
measure prevention where a pre-settlement gate exists. In the delegated flow the
predicate guard is evaluated \emph{before} the facilitator is called, and
settlement proceeds only if the guard allows it. Across nine scenarios
(payee substitution, revoked merchant, amount escalation, asset switch, chain
switch, task substitution, expired authorization, selection-then-substitution,
and escalation-then-replay) the guard denied eight outright and the facilitator
rejected the ninth on an already-consumed nonce; no settlement receipt was
issued for any attacked attempt (Table~\ref{tab:presettlement}).

\begin{table}[t]
\centering
\small
\caption{Pre-settlement prevention in the delegated flow. ``Receipt issued''
refers to the attacked attempt; the guard is queried before any settlement call,
so a false value is direct evidence that no value moved.}
\label{tab:presettlement}
\begin{tabular}{@{}lccc@{}}
\toprule
Scenario & Guard denied & Settlement attempted & Receipt issued \\
\midrule
Malicious merchant: payee substitution      & yes & no & no \\
Malicious merchant: revoked identity        & yes & no & no \\
Compromised agent: amount escalation        & yes & no & no \\
Compromised agent: asset switch             & yes & no & no \\
Compromised agent: chain switch             & yes & no & no \\
Compromised agent: task substitution        & yes & no & no \\
Expired authorization                       & yes & no & no \\
Multi-step: selection then substitution     & yes & no & no \\
Multi-step: escalation then replay          & n/a & yes & no \\
\bottomrule
\end{tabular}
\end{table}

\paragraph{Concurrent transactions.}
Concurrency is the remaining failure mode a sequential replay test cannot
observe: two requests can both pass a consumed-nonce check before either writes.
We therefore reserve the nonce in an in-flight set between signature
verification and settlement, and evaluate the reservation under burst load at
concurrency levels 2, 4, 8, 16 and 32, with twenty trials per level in three
shapes: identical nonces, distinct nonces, and a mixture. Across 3{,}720
concurrent attempts no nonce settled twice, every distinct-nonce burst settled
completely, and no anomaly occurred (Table~\ref{tab:concurrency}).

\begin{table}[t]
\centering
\small
\caption{Concurrent settlement. Each cell aggregates twenty trials; the
same-nonce column is the number of settlements observed, which must equal the
trial count (one per trial).}
\label{tab:concurrency}
\begin{tabular}{@{}rrrrr@{}}
\toprule
Concurrency & Trials per shape & Settled (same nonce) & Settled (distinct) & Anomalies \\
\midrule
2  & 20 & 20  & 40  & 0 \\
4  & 20 & 20  & 80  & 0 \\
8  & 20 & 20  & 160 & 0 \\
16 & 20 & 20  & 320 & 0 \\
32 & 20 & 20  & 640 & 0 \\
\bottomrule
\end{tabular}
\end{table}

\paragraph{Scope.}
Three limits apply. First, the matrix measures detection; prevention is
established only for the delegated flow, where a guard exists before settlement
--- in the exact flow the human signs the concrete descriptor, so there is no
pre-settlement gate to test beyond signature validity. Second, the nonce
reservation is per-process; a multi-replica facilitator requires a shared nonce
store, which we do not claim to have evaluated. Third, the primitives are
modelled from the protocol surface rather than discovered by fuzzing, so the
matrix bounds the attacks we thought of, not all attacks.
